\documentclass[11pt]{article}

% ---------------------------------------------------------------- packages
\usepackage[T1]{fontenc}
\usepackage[utf8]{inputenc}
\usepackage{libertine}
\usepackage[varqu,varl,scaled=0.95]{zi4}
\usepackage{amsmath}
\usepackage[margin=1in,headheight=14pt]{geometry}
\usepackage{microtype}
\usepackage{xcolor}
\usepackage{booktabs}
\usepackage{tabularx}
\usepackage{xltabular}
\usepackage{array}
\usepackage{enumitem}
\usepackage{listings}
\usepackage[most]{tcolorbox}
\usepackage{titlesec}
\usepackage{fancyhdr}
\usepackage[hidelinks]{hyperref}
\usepackage{xspace}
\usepackage{float}

% ---------------------------------------------------------------- colors
\definecolor{ink}{HTML}{1B1F24}
\definecolor{accent}{HTML}{1F4E79}
\definecolor{finding}{HTML}{9C3D10}
\definecolor{findingbg}{HTML}{FBF1EB}
\definecolor{proposal}{HTML}{1B6B5A}
\definecolor{proposalbg}{HTML}{ECF6F3}
\definecolor{codebg}{HTML}{F5F6F8}
\definecolor{rule}{HTML}{C9CED6}
\definecolor{good}{HTML}{1B6B3A}
\definecolor{part}{HTML}{8A6100}
\definecolor{bad}{HTML}{A8261B}
\definecolor{comment}{HTML}{6A737D}
\definecolor{kw}{HTML}{6F2DA8}
\definecolor{str}{HTML}{0B6B3A}

\color{ink}
\hypersetup{colorlinks=true,linkcolor=accent,urlcolor=accent,citecolor=accent}

% ---------------------------------------------------------------- headings
\titleformat{\section}{\Large\bfseries\color{accent}}{\thesection}{0.8em}{}
\titleformat{\subsection}{\large\bfseries\color{accent}}{\thesubsection}{0.7em}{}
\titleformat{\subsubsection}{\normalsize\bfseries}{\thesubsubsection}{0.6em}{}
\titlespacing*{\section}{0pt}{2.2ex plus 1ex}{1.0ex}
\titlespacing*{\subsection}{0pt}{1.8ex plus 0.8ex}{0.7ex}

\pagestyle{fancy}
\fancyhf{}
\fancyhead[L]{\small\color{comment}Who Reaches for Combinatorial Testing}
\fancyhead[R]{\small\color{comment}UnitTestDesign.jl review}
\fancyfoot[C]{\small\color{comment}\thepage}
\renewcommand{\headrulewidth}{0.3pt}

\setlist{itemsep=0.25ex,topsep=0.6ex}
\setlength{\parskip}{0.55ex}

% ---------------------------------------------------------------- code
\lstdefinelanguage{Julia}{
  morekeywords={abstract,break,catch,const,continue,do,else,elseif,end,export,
    false,for,function,if,import,in,let,macro,module,mutable,nothing,quote,
    return,struct,true,try,using,while,begin},
  sensitive=true,
  morecomment=[l]\#,
  morestring=[b]",
}
\lstset{
  language=Julia,
  basicstyle=\ttfamily\small,
  keywordstyle=\color{kw},
  commentstyle=\color{comment},
  stringstyle=\color{str},
  columns=fullflexible,
  keepspaces=true,
  showstringspaces=false,
  upquote=true,
  breaklines=true,
  aboveskip=0.6ex,
  belowskip=0.2ex,
}
\newtcblisting{jlcode}[1][]{
  listing only, breakable, enhanced,
  colback=codebg, colframe=codebg, boxrule=0pt, arc=2pt,
  left=6pt, right=6pt, top=3pt, bottom=3pt,
  listing options={language=Julia,basicstyle=\ttfamily\small},
  #1
}
\newtcblisting{term}[1][]{
  listing only, breakable, enhanced,
  colback=white, colframe=rule, boxrule=0.5pt, arc=2pt,
  left=6pt, right=6pt, top=3pt, bottom=3pt,
  listing options={language={},basicstyle=\ttfamily\footnotesize,columns=fixed,keepspaces=true},
  #1
}

% ---------------------------------------------------------------- boxes
\newtcolorbox{findingbox}[1]{
  breakable, enhanced, colback=findingbg, colframe=finding, boxrule=0pt,
  leftrule=3pt, arc=0pt, left=8pt, right=8pt, top=4pt, bottom=4pt,
  fonttitle=\bfseries\color{finding}, coltitle=finding, colbacktitle=findingbg,
  title={#1}, attach title to upper={\ \ }
}
\newtcolorbox{proposalbox}[1]{
  breakable, enhanced, colback=proposalbg, colframe=proposal, boxrule=0pt,
  leftrule=3pt, arc=0pt, left=8pt, right=8pt, top=4pt, bottom=4pt,
  fonttitle=\bfseries\color{proposal}, coltitle=proposal, colbacktitle=proposalbg,
  title={#1}, attach title to upper={\ \ }
}

% ---------------------------------------------------------------- shorthands
\newcommand{\code}[1]{\texttt{\detokenize{#1}}}
\newcommand{\utd}{\textsf{UnitTestDesign.jl}\xspace}
\newcommand{\met}{\textcolor{good}{\textbf{Met}}}
\newcommand{\partly}{\textcolor{part}{\textbf{Partly}}}
\newcommand{\notyet}{\textcolor{bad}{\textbf{Not yet}}}
\newcolumntype{L}[1]{>{\raggedright\arraybackslash}p{#1}}
\newcolumntype{Y}{>{\raggedright\arraybackslash}X}
\renewcommand{\arraystretch}{1.18}

% =====================================================================
\begin{document}

\begin{titlepage}
\setlength{\parindent}{0pt}
\vspace*{1.2in}
{\color{accent}\rule{\linewidth}{1.2pt}}\\[1.2ex]
{\fontsize{24}{29}\selectfont\bfseries\color{accent} Who Reaches for Combinatorial Testing,\\[0.3ex] and What They Need From It}\\[2ex]
{\Large A user-centered review of \utd, with experiments and a roadmap}\\[1.2ex]
{\color{accent}\rule{\linewidth}{0.6pt}}

\vspace{2.2em}
\begin{tabular}{@{}ll@{}}
Prepared for & Andrew Dolgert \\
Prepared by  & Claude (Opus 5.5), working in the repository \\
Date         & September 23, 2026 \\
Reviewed     & \utd v0.4.0, commit \code{e0504d6}, branch \code{feature/upgrade} \\
Environment  & Julia 1.12.5 on macOS \\
\end{tabular}

\vspace{3em}
\begin{minipage}{0.9\linewidth}
\small\color{comment}
Statements marked \emph{measured} come from experiments I ran during this
review; Appendix~\ref{app:repro} lists them. Function names in the proposal
sections (\code{TestSpace}, \code{coverage}, \code{diagnose}, and so on) are
sketches of a possible interface, not existing API. I did not read the
earlier \code{interface_*} analyses in the repository.
\end{minipage}
\vfill
\end{titlepage}

{\small\setlength{\parskip}{0pt}\tableofcontents}
\newpage

% =====================================================================
\section*{Summary}
\addcontentsline{toc}{section}{Summary}

\utd answers one question well: \emph{given a list of values for each
parameter, which few combinations should I test so that every pair (or
triple, or $t$-tuple) of values appears somewhere?} Its default engine, IPOG, is
fast. Fifteen four-valued parameters at strength four take about four seconds
and produce 958 tests, within 4\% of the 924 reported in the IPOG paper~\cite{lei2008}. In each
of the 888 constrained random problems where IPOG returned a result, an
independent checker confirmed that the result was valid and complete. The core
algorithms are sound wherever they finish.

The package's limits lie around generation rather than in it. People reach
for combinatorial testing at a particular moment: a nested loop over options
has just exploded, or each run costs an hour of cluster time, or a bug arrived
from a combination nobody tried. At that moment they have questions. How many
tests will this be? Did it really cover everything? What did it leave out,
and why? When test 7 fails, which combination is to blame? What would it cost
to go further? Today the package returns a \code{Vector{Vector{Any}}} and says
nothing more. Users have to work out everything else themselves.

What I found, in order of consequence:

\begin{enumerate}[leftmargin=1.5em]
\item \textbf{Constraints are where users get hurt.} With a \code{disallow}
  rule, IPOG crashed with \code{BoundsError: attempt to access ... at index [0]}
  on 8.6\% of random pairwise problems and 13.8\% of three-way problems
  (\emph{measured}). Every problem that contained an \emph{implicit}
  constraint crashed. An implicit constraint arises when two innocent rules
  together make some pair impossible. On the same kind of input GND does
  not crash; it loops forever. And because unassigned parameters reach the rule
  as \code{nothing}, a parameter whose real values include \code{nothing} (a
  common Julia default) silently loses required coverage.
\item \textbf{GND has a one-token bug.} One comparison uses \code{n_way} where
  it needs \code{n_way - 1}. As a result, once as many parameters are set as
  the strength (two, for pairwise), GND picks every remaining value at random.
  Fixing it shrinks GND's designs by 17--35\%, and GND then beats IPOG when parameters have eight values
  (\emph{measured}). The documentation describes GND as the engine that finds
  shorter designs, but the shipped code finds longer ones.
\item \textbf{Parameters have positions but no names.} That single decision
  makes constraints positional, makes forty-parameter configurations
  unreadable, and makes a failing case print as
  \code{Any[1000, :newton, 0.001, true]}.
\item \textbf{Several natural users are not served at all.} These include the
  maintainer who wants to know what their existing hand-written tests cover,
  the developer debugging a failure, the engineer sizing a CI matrix, the
  tester of invalid inputs, and the AI agent writing tests on someone's behalf.
\end{enumerate}

The recommendation, in one sentence: \textbf{turn the package from a
generator that returns a list into a short conversation about a test space}.
That conversation runs: plan, generate with a stated guarantee, run with
names, diagnose, extend. Each step reports what it did and what is possible
next. Table~\ref{tab:top} lists the ten changes I would make first.

\begin{table}[H]
\centering\small
\caption{The ten changes I would make first.}\label{tab:top}
\begin{tabularx}{\linewidth}{@{}r L{0.50\linewidth} L{0.11\linewidth} Y@{}}
\toprule
 & Change & Effort & Who benefits \\
\midrule
1 & Fix GND's match condition (one token); seed GND by default and record the seed. & Hours & Everyone using GND \\
2 & Make constraints safe: drop and report impossible tuples, keep every row completable, never loop without progress. & Days & Everyone with rules \\
3 & Return a design object that prints its guarantee, exclusions, and next steps, while still iterating like today's vector. & Days & Everyone; reviewers; agents \\
4 & Named parameters; each case a \code{NamedTuple}. & Days & Everyone \\
5 & Constraints scoped by argument names, so rules never see \code{nothing}. & Days & Everyone with rules \\
6 & \code{coverage(cases, space)} to measure any test set, including hand-written ones; \code{topup} to fill the gaps. & Days & Maintainers, auditors \\
7 & \code{run_cases} and \code{diagnose}: name cases in test output, rank suspect combinations, propose follow-ups. & 1--2 weeks & Anyone with a failure \\
8 & Error messages in the user's vocabulary (parameter names and values), checked before generation. & Days & Everyone; especially agents \\
9 & Export to Tables.jl, CSV/TOML, a GitHub Actions matrix, and literal test code. & Days each & CI, simulation, agents \\
10 & Organize docs by job, run examples as doctests, correct the engine guidance, add a one-page agent guide. & Days & First-time users \\
\bottomrule
\end{tabularx}
\end{table}

% =====================================================================
\section{How I did this review}

I read all of the source, tests, and documentation, plus the JuliaCon paper.
Then I used the package the way a newcomer would, in a scratch environment,
and ran experiments whose results appear throughout this report:

\begin{itemize}[leftmargin=1.5em]
\item An \textbf{independent coverage checker} enumerates every feasible
  $t$-tuple by brute force and shares no code with the package. A tuple is
  \emph{feasible} if at least one complete, allowed test contains it.
\item \textbf{1,000 random constrained problems} (500 at strength 2 and 500
  at strength 3). Each had 4--6 parameters with 2--4 values, and 1--3 forbidden
  partial assignments over 2--3 parameters. A fixed seed makes the run
  repeatable.
\item \textbf{Other checks:} timing and allocation across problem sizes, GND
  under a watchdog, verification of mixed strength, stability under small
  edits, and the printed output of failing tests.
\item A \textbf{throwaway prototype} of about 50 lines, to check that the
  constraint fix I recommend in Section~\ref{sec:solvers} actually works.
\end{itemize}

The scripts are small enough to become regression tests.

% =====================================================================
\section{The moment someone reaches for combinatorial testing}\label{sec:moment}

\subsection{The situation it fits}

Combinatorial testing fits when three conditions hold at once.

\begin{enumerate}[leftmargin=1.5em]
\item \textbf{The inputs are, or can be made into, discrete choices:}
  options, modes, flags, algorithms, element types, container types,
  platforms, or representative values of equivalence classes.
\item \textbf{Failures come from interactions among a few choices.} Kuhn,
  Wallace, and Gallo found that the failures they studied were triggered by
  at most six interacting parameters, and most by one or
  two~\cite{kuhn2004}. Covering every $t$-way combination therefore comes
  close to exhaustive testing at a small fraction of the cost.
\item \textbf{Exhaustive testing is unaffordable.} Either the product is huge
  (twenty four-valued options give about $10^{12}$ combinations), or each test is
  expensive (a simulation that runs for an hour).
\end{enumerate}

\subsection{The moments that trigger it}

Nobody sets out to do combinatorial testing. People arrive at it from a few
recognizable situations, and the package's first page should name them.

\begin{description}[leftmargin=1.5em,style=nextline]
\item[The loop that exploded.] Someone writes
  \code{for T in types, A in arrays, alg in algorithms, ...}, and the suite goes
  from seconds to an hour. Or they quietly comment out half the loop.
\item[The expensive run.] Each case is a cluster job. The budget is thirty
  runs, not three thousand.
\item[The bug from nowhere.] A user reports a crash with a combination of
  options nobody tried, and the maintainer wonders which other combinations
  have never run.
\item[The matrix bill.] The CI matrix of operating systems, Julia versions,
  architectures, thread counts, and dependency versions costs too much. Or it
  hits GitHub Actions' limit of 256 jobs per matrix~\cite{ghmatrix}.
\item[The audit.] A reviewer, funder, or regulator asks how thoroughly the
  options have been exercised.
\end{description}

Generation serves the first four situations. The audit needs
\emph{measurement}, which the package does not yet offer. And every situation
eventually leads to a failure that needs \emph{diagnosis}, which it also
doesn't offer.

\subsection{Where it does not fit}

Part of serving users is telling them when another tool is better. The
documentation says little about this today.

\begin{itemize}[leftmargin=1.5em]
\item \textbf{Finding the one value that breaks a numerical routine} calls for
  property-based testing or fuzzing, for example
  Supposition.jl~\cite{supposition}. It explores values within domains and
  shrinks failures to minimal examples.
\item \textbf{Estimating how much each factor affects an outcome} is design of
  experiments. Orthogonal arrays and fractional factorials are balanced for
  estimation, while covering arrays are not. A scientist planning a
  sensitivity study should be pointed elsewhere.
\item \textbf{Sweeps over continuous parameters} call for space-filling
  samples such as Sobol sequences or Latin hypercubes.
\end{itemize}

Property-based testing is a partner, not a rival. Combinatorial testing
decides \emph{which partitions to combine}; property-based testing decides
\emph{which values to draw within each partition}. The package's Testing
Methods page already hints at this, and it could become a first-class feature
(Section~\ref{sec:api}).

% =====================================================================
\section{Who the users are}\label{sec:users}

Table~\ref{tab:personas} lists eight kinds of users. The first four are
implied by the current documentation. The last four are not addressed at
all, though each has a well-established counterpart in the testing
literature or in other tools.

\begin{table}[h]
\centering\small
\caption{Who reaches for this package, and whether it serves them today.}\label{tab:personas}
\begin{tabularx}{\linewidth}{@{}L{0.20\linewidth} Y L{0.14\linewidth}@{}}
\toprule
User & What they need most & Served today \\
\midrule
Generic-code author & Combinations of types and options; readable failures; fast generation inside \code{runtests.jl} & \partly \\
Simulation developer & Few expensive runs; constraints; reference runs included; cases as files for a scheduler & \partly \\
CI / release engineer & Matrix output; rules such as ``no arm64 on Windows''; the same jobs from commit to commit & \notyet \\
Test-data builder & Tables with named columns that cover value combinations & \partly \\
\midrule
Maintainer / auditor & Measure the interaction coverage of existing tests; add only what is missing & \notyet \\
Debugger & Which combination caused this failure, and what to run next & \notyet \\
Robustness tester & Invalid values tested one at a time, so one does not mask another & \notyet \\
AI coding agent & Knowing when to use it; results that describe themselves; checkable claims; actionable errors & \notyet \\
\bottomrule
\end{tabularx}
\end{table}

\subsection{Users the documentation already imagines}

\paragraph{The generic-code author.}
This user is especially common in Julia. Multiple dispatch means that
combinations of element type, container type, dimension, and option really
are different compiled methods. The extended Hilbert-curve example in the
documentation is this user's story, and combinatorial testing found a real
bug there. They run generation at test time, so it must be quick and
deterministic. They need failures they can read. They care about stability,
so that a failure on Tuesday can be compared with Wednesday's. Positional,
unnamed parameters and \code{Vector{Any}} results work against all of these.

\paragraph{The simulation developer.}
This is the user in the paper's statement of need. They have a configuration
file with dozens of options, and each run takes minutes to hours, often on a
cluster. Their ``tests'' are rarely \code{@test} calls. They are jobs, so they
need the cases as data (TOML, CSV, JSON) to hand to a scheduler. Their options
come with many rules, since some settings only make sense together, and that
is exactly where the package is most fragile today. They have reference runs
that must be included. They want stronger coverage on risky subsystems, which
the \code{wayness} option provides. When a campaign is half done and some runs
crashed, they need to know what remains untested, and the package cannot tell
them. Forty positional parameters are also hard to maintain.

\paragraph{The CI and release engineer.}
Their space is operating system $\times$ Julia version $\times$ architecture
$\times$ threads $\times$ dependency versions. They need the design as a
GitHub Actions \code{include:} list and rules that express what is
impossible. They also need stability: if adding a Julia version reshuffles
every job, yesterday's failure can no longer be compared with today's. There
is also an opportunity. A randomized engine seeded with the run number
would give a different pairwise matrix on each CI run, so higher-order
coverage accumulates across runs at no extra cost.

\paragraph{The test-data builder.}
The documentation's DataFrame example builds rows that cover value
combinations for testing a data pipeline. Today the user has to convert
positional vectors into named columns by hand.

\subsection{Users no one has addressed}

\paragraph{The maintainer or auditor of an existing suite.}
Most projects already have tests. The first question a thoughtful maintainer
asks is not ``generate new tests'' but ``what do my current tests cover, and
what is the smallest set I could add?'' NIST built a tool, CCM, for exactly
this question~\cite{ccm}. \utd already has the internal machinery: coverage
structures, and seeds that count toward coverage. It exposes none of it. A
\code{coverage} function and a \code{topup} function would serve this user
without disturbing how they already test, and this is the gentlest way to
bring combinatorial testing into a project.

\paragraph{The debugger.}
After a failure, the natural question is \emph{which combination caused
it?} A covering design is unusually good at answering this. Candidate causes
are combinations that appear in failing cases but never in passing ones, and
a few follow-up tests can separate them. Tools such as BEN build on this
idea~\cite{ben}. There is a second point, known as \emph{masking}. Tuples
that appeared only in failing cases have not been shown to work, so the
coverage guarantee is weaker than it looks until they are re-tested
elsewhere~\cite{yilmaz2014}. Section~\ref{sec:ev-failure} shows what such a
report could say for a real run.

\paragraph{The robustness tester.}
Invalid inputs (negative sizes, empty strings, \code{NaN} tolerances) should
appear \emph{one per test case}. Otherwise the first invalid value triggers
an error and hides the others. Microsoft's PICT supports this by letting
users mark values as invalid~\cite{pict}. It is a small addition that
changes what the package is good for.

\paragraph{The AI coding agent.}
Agents now write a large share of tests. Left alone, they tend either to
write the Cartesian loop that explodes or to hand-pick three cases, and this
package is the principled middle path. An agent needs the package to
(i)~state in its docstrings the situation it is for, so the agent knows to
reach for it; (ii)~return results that describe themselves in text, so the
agent can report accurately; (iii)~be deterministic, so reruns match;
(iv)~produce errors that name the parameter and suggest a fix, because agents
act on messages literally; (v)~offer a \code{verify} or \code{coverage}
function, so claims can be checked rather than trusted; and (vi)~optionally
emit the design as literal test code, so a human reviewer sees every case in
the diff. The human who reviews an agent's tests needs (ii) and (v) just as
much.

% =====================================================================
\section{What users worry about, and how the package does today}\label{sec:concerns}

Whatever brings them to the package, users share a common set of concerns.
Table~\ref{tab:concerns} lists each one with an assessment and a pointer to
the evidence.

{\small
\begin{xltabular}{\linewidth}{@{}L{0.25\linewidth} L{0.08\linewidth} Y@{}}
\caption{Users' concerns, rated against version 0.4.0.}\label{tab:concerns}\\
\toprule
The user's question & Today & Evidence \\
\midrule
\endfirsthead
\toprule
The user's question & Today & Evidence \\
\midrule
\endhead
\bottomrule
\endlastfoot
Is everything really covered? & \partly & Whenever a design is returned, the guarantee holds: 888 of 888 constrained designs, and mixed strength (\S\ref{sec:ev-mixed}). But nothing lets the user see or check that; no coverage function is exported. \\
How many tests will this be, and is that good? & \notyet & There is no preview, no comparison with the full factorial, and no lower bound. \code{full_factorial} allocates before checking size: sixteen four-valued parameters would need 550\,GB. \\
Will it respect what my function accepts? & \notyet & IPOG crashes on 8.6\% and 13.8\% of constrained problems; GND hangs; natural rule code throws; a real \code{nothing} value silently loses coverage (\S\ref{sec:ev-constraints}). \\
Will my must-run cases be included? & \partly & Yes. But a typo in a value produces \code{MethodError: Cannot convert ... Nothing to ... Int64}; must-run cases that break the rules are kept silently; a forbidden excursion base case silently disappears. \\
Will I get the same tests tomorrow? & \partly & IPOG is deterministic. GND is unseeded by default, so two calls differ. Adding one value to one parameter kept only 3 of 10 IPOG cases (\S\ref{sec:ev-stability}). \\
When a case fails, which one was it? & \notyet & The README loop reports ``2 errors'' with no parameters. The per-case pattern prints \code{Any[1000, :newton, 0.001, true]} (\S\ref{sec:ev-failure}). \\
Which combination caused it? & \notyet & Nothing is provided, although the data to answer it is already in hand. \\
Can I fit this in my budget? & \partly & IPOG puts most coverage early (16 of 28 cases reach 90\% of pairs), but this isn't documented, and there is no \code{max_cases} (\S\ref{sec:ev-curve}). \\
Can I test the risky parts harder? & \partly & Mixed strength works and is verified. But the documented example throws, the user's \code{Dict} is modified, and GND fails an assertion in one case (\S\ref{sec:ev-mixed}). \\
Can I test invalid inputs properly? & \notyet & There is no notion of invalid values. \\
Which engine should I use? & \notyet & The documentation says GND is shorter. It is 26--38\% longer, because of a bug (\S\ref{sec:ev-gnd}). \\
Is it light and fast? & \met & Two dependencies. IPOG runs in well under a second for typical sizes (\S\ref{sec:ev-perf}). \\
Will it fit my tools? & \partly & It works inside a \code{for} loop. It has no names, no Tables.jl support, no file or matrix export, and no Test.jl integration. \\
\end{xltabular}
}

% =====================================================================
\section{Evidence}\label{sec:evidence}

\subsection{Constraints}\label{sec:ev-constraints}

\subsubsection*{The \code{nothing} convention}

Rules are ordinary Julia functions over all parameters. During generation,
the package calls a rule on partial test cases and passes \code{nothing} for
parameters it has not chosen yet (\code{src/factorial_interface.jl:85}). The
guide warns about this, but natural rule code still breaks:

\begin{jlcode}
all_pairs([1,2,3], [0.0, 1.0, 10.0], [:a,:b];
          disallow = (n, x, s) -> n > 2 && x > 1.0)
# MethodError: no method matching isless(::Float64, ::Nothing)
\end{jlcode}

The convention has a subtler cost as well. In Julia, \code{nothing} is a
common \emph{real} value, as in \code{maxiter = nothing}. When a parameter
takes \code{nothing} as one of its values, the package cannot tell ``not yet
chosen'' from ``chosen to be \code{nothing}''. It then drops required tuples
without saying so:

\begin{jlcode}
maxiter = [nothing, 10, 100]; mode = [:fast, :slow]
verbose = [true, false];      scale = [1.0, 2.0, 3.0]
dis = (m, md, v, s) -> m === nothing && md === :slow
cases = all_pairs(maxiter, mode, verbose, scale; disallow = dis)
# Independent check: 35 of 36 feasible pairs covered.
# Never tested: mode = :slow, scale = 1.0   (a perfectly valid pair)
\end{jlcode}

\subsubsection*{Random constrained problems}

Table~\ref{tab:random} summarizes 1,000 random problems. IPOG never
returned a wrong answer. It either returned a correct one or crashed. For
every crash I checked whether some tuple that the rule permits in isolation
is impossible in any complete, allowed test, which is the definition of an
implicit constraint. The rates depend on how tightly the rules bind, and
these small, fairly constrained problems are not a sample of real projects.
The failure classes are real and easy to hit, as the next example shows.

\begin{table}[h]
\centering\small
\caption{IPOG on 1,000 random constrained problems (\emph{measured}).}\label{tab:random}
\begin{tabular}{@{}lrr@{}}
\toprule
 & Strength 2 & Strength 3 \\
\midrule
Problems & 500 & 500 \\
Returned; independently verified valid and complete & 457 & 431 \\
Returned but invalid or incomplete & 0 & 0 \\
Crashed with \code{BoundsError} & 43 (8.6\%) & 69 (13.8\%) \\
\quad of which contained an implicit constraint & 28 & 65 \\
\quad of which were greedy dead ends (no implicit constraint) & 15 & 4 \\
Problems with an implicit constraint that did \emph{not} crash & 0 & 0 \\
\midrule
Prototype with completion checks (\S\ref{sec:solvers}): valid and complete & 500 & 500 \\
\bottomrule
\end{tabular}
\end{table}

Implicit constraints are ordinary in practice. Here is a three-parameter
example:

\begin{jlcode}
os = [:linux, :mac, :windows]; gpu = [false, true]; driver = [:cuda, :none]
dis(o, g, d) = (g === true && d === :none) || (o === :windows && d === :cuda)
all_pairs(os, gpu, driver; disallow = dis)
# BoundsError: attempt to access 2-element Vector{Symbol} at index [0]
all_pairs(os, gpu, driver; disallow = dis, engine = GND())
# never returns (killed by a 45-second watchdog; 12 combinations in total)
\end{jlcode}

Neither rule mentions Windows together with a GPU. But a GPU needs CUDA and
Windows forbids CUDA, so no valid test can pair \code{os = :windows} with
\code{gpu = true}. The user sees an index error that points into the
package's internals, or a process that never finishes.

\begin{findingbox}{Why it happens.}
IPOG keeps any tuple that the rule does not reject on its own. When the tuple
fits nowhere, IPOG appends it as a new row (\code{src/parameter_order.jl:170}).
The final fill then finds no allowed value for some slot and leaves a zero
(lines 224--226). Converting back to values then indexes \code{p[0]}. GND's
outer loop (\code{src/greedy_tuples.jl:195}) waits forever for a tuple no
candidate can cover. Its inner loop (line 201) can also spin if every
candidate is forbidden. The greedy dead ends have a related cause: values
are placed after checking only that the row is not \emph{yet} forbidden, not
that it can still be completed. The test suite exercises only simple rules
over two parameters, so neither failure is caught.
\end{findingbox}

\subsubsection*{Silent acceptance}

Two more behaviors pass without a word. Must-run cases (\code{seeds}) that
violate the rules are included in the output anyway. And when the base case
of an excursion is forbidden, it is dropped, even though the base case is
the whole point of an excursion.

\subsection{GND picks most values at random}\label{sec:ev-gnd}

\code{most_matches_existing} (\code{src/coverage_matrix.jl:193}) scores each
candidate value by how many uncovered tuples it would complete. The
comment above it says a tuple can match at most $n-1$ known values, because
the parameter being chosen is still unset. The code, however, tests
\code{match_cnt == min(max_known, n_way)}. Once the number of set parameters
reaches the strength (two, for pairwise), no tuple satisfies this, every score is zero, and the tie-breaker picks a
random value. GND then relies on its \code{M} random restarts to find
anything good.

\begin{table}[h]
\centering\small
\caption{Number of test cases: IPOG, GND as shipped, and GND with \code{n_way - 1} (three seeds each; \emph{measured}).}\label{tab:gnd}
\begin{tabular}{@{}lrrr@{}}
\toprule
Parameters $\times$ values, strength & IPOG & GND as shipped & GND patched \\
\midrule
$10\times4$, $t=2$ & 28 & 36--38 & 30 \\
$10\times8$, $t=2$ & 113 & 151--153 & \textbf{100} \\
$20\times4$, $t=2$ & 37 & 48--51 & 37 \\
$10\times4$, $t=3$ & 143 & 180--186 & 144--150 \\
\bottomrule
\end{tabular}
\end{table}

Increasing \code{M} does not close the gap. On the $10\times4$ problem, GND
as shipped gave 40--41 cases at \code{M=10}, 36--38 at 50, 32--34 at 200, and
32--33 at 1000 (1.2\,s). Even twenty times the default effort did not reach
IPOG's 28. The patched version still covered all 720 pairs when checked
independently.

\subsection{Mixed strength works, but its interface is fragile}\label{sec:ev-mixed}

\code{all_pairs(fill(1:4, 10)...; wayness = Dict(3 => [[3,4,5,6]]))} produced
64 cases (IPOG) and 84 (GND). The independent check confirmed all 720 pairs
and all 256 triples on parameters 3--6. The feature itself is sound. Around
it:

\begin{itemize}[leftmargin=1.5em]
\item The guide's own example, \code{Dict(3 => [[3:6], [25:30]])}, throws
  \code{MethodError: Cannot convert an object of type Int64 to an object of
  type UnitRange{Int64}}.
\item The user's dictionary is modified. After a call it has gained an entry
  \code{2 => [[1, 2, 3, 4]]} (\code{src/parameter_order.jl:436},
  \code{src/excursions.jl:32}), so reusing it in a later call changes that
  call's meaning.
\item GND fails with \code{AssertionError: minimum(orders) > base_wayness} if
  any listed strength equals the base strength
  (\code{src/coverage_matrix.jl:426}).
\end{itemize}

\subsection{Performance}\label{sec:ev-perf}

IPOG is fast enough that generation can run inside a test suite
(Table~\ref{tab:perf}). Allocation is heavy because hot loops slice matrix
columns, but it isn't a practical problem at these sizes. The one real risk
is \code{full_factorial}, which allocates its whole result before checking
whether the size is reasonable.

\begin{table}[H]
\centering\small
\caption{Generation time and total allocation (\emph{measured}; times include some compilation; allocation is cumulative, not peak).}\label{tab:perf}
\begin{tabular}{@{}lrrrrr@{}}
\toprule
Problem & IPOG cases & IPOG time & Allocated & GND cases & GND time \\
\midrule
$10\times4$, $t=2$  & 28  & 0.12\,s & 42\,MB  & 36--38 & 0.08\,s \\
$20\times4$, $t=2$  & 37  & 0.11\,s & 42\,MB  & 48--51 & 0.21\,s \\
$50\times5$, $t=2$  & 75  & 0.24\,s & 0.4\,GB & -- & -- \\
$100\times3$, $t=2$ & 33  & 0.37\,s & 0.5\,GB & -- & -- \\
$10\times4$, $t=3$  & 143 & 0.11\,s & 87\,MB  & 180--186 & 1.02\,s \\
$30\times4$, $t=3$  & 256 & 1.55\,s & 5.7\,GB & -- & -- \\
$50\times3$, $t=3$  & 136 & 1.87\,s & 8.6\,GB & -- & -- \\
$15\times4$, $t=4$  & 958 & 4.06\,s & 17\,GB  & -- & -- \\
\bottomrule
\end{tabular}
\end{table}

\subsection{What a failure looks like}\label{sec:ev-failure}

Here is a solver with a two-way bug. It fails only when
\code{method = :newton} and \code{sparse = true}.

\begin{jlcode}
solve(n, method, tol, sparse) =
    (method == :newton && sparse) ? error("singular factorization") : n * tol
cases = all_pairs([10, 100, 1000], [:newton, :bicg, :gmres], [1e-3, 1e-6], [false, true])
\end{jlcode}

With the README's loop, Test.jl reports \code{solver | Pass 8 Error 2} and two
stack traces for the expression \code{solve(tc...) > 0}. The output never
says which parameters failed. With the documentation's per-case pattern
(\code{@testset "solver $tc" for tc in cases}), the failing case prints as
\code{solver Any[1000, :newton, 0.001, true]}, which a reader has to decode by
position.

The design and the outcomes already contain the answer. About twenty lines of
code, run on this very output, find the combinations that appear in failing
cases but never in passing ones:

\begin{term}
suspicious: method=:newton, sparse=true   (in 2 of 2 failures, 0 of 8 passes)
suspicious: n=1000, method=:newton        (in 1 of 2 failures, 0 of 8 passes)
suspicious: n=100, sparse=true            (in 1 of 2 failures, 0 of 8 passes)
\end{term}

The true cause ranks first. The other two pairs are \emph{masked}: they
appeared only in failing cases, so nobody knows whether they work. Two
follow-up cases would settle all three:
\code{(1000, :newton, 0.001, false)} and \code{(100, :gmres, 1e-6, true)}.
Section~\ref{sec:communicate} turns this into a proposed feature.

\subsection{Reproducibility and stability}\label{sec:ev-stability}

\code{GND()} seeds itself from the global \code{MersenneTwister()}
(\code{src/factorial_interface.jl:61}). Two calls with identical arguments
gave different designs, so a failure found with GND cannot be reproduced
unless the user thought to pass an RNG in advance. IPOG is deterministic but
not stable: adding one value to one of five parameters changed the design
from 10 to 14 cases, and only 3 of the original rows survived. This matters
to anyone who compares failures across commits. Passing the old design as
\code{seeds} would preserve it, but nothing suggests doing so.

\subsection{Coverage accumulates early, and higher coverage comes free}\label{sec:ev-curve}

The order of the cases matters for anyone on a budget. On $10\times4$
pairwise, IPOG's first cases cover pairs quickly:

\begin{center}\small
\begin{tabular}{@{}lrrrrrrr@{}}
\toprule
First $j$ cases & 4 & 8 & 12 & 16 & 20 & 24 & all \\
\midrule
IPOG, share of pairs covered & 25\% & 47\% & 68\% & 90\% & 93\% & 97\% & 100\% (28) \\
GND as shipped & 25\% & 47\% & 64\% & 77\% & 86\% & 92\% & 100\% (37) \\
\bottomrule
\end{tabular}
\end{center}

A pairwise design also covers a good share of triples as a side effect: 25\%
for $5\times4$, 40\% for $10\times4$, 54\% for $10\times3$, and 61\% for
$20\times3$ (\emph{measured}). Users never see either number, yet both
answer the questions ``can I stop early?'' and ``how much am I getting?''

\subsection{Error messages}\label{sec:ev-errors}

Most error messages come from deep inside the algorithms and use the
algorithms' vocabulary, not the user's (Table~\ref{tab:errors}).

{\small
\begin{xltabular}{\linewidth}{@{}L{0.25\linewidth} L{0.33\linewidth} Y@{}}
\caption{What users see, and what would help (\emph{measured}).}\label{tab:errors}\\
\toprule
Input & What the user sees & What would help \\
\midrule
\endfirsthead
\toprule
Input & What the user sees & What would help \\
\midrule
\endhead
\bottomrule
\endlastfoot
A parameter with one value, \code{[true]} & \code{DomainError with Bool[1]: Each argument should be a list of parameter values} & Accept it; a fixed setting is legitimate. \\
Strength 3 with two parameters & \code{BoundsError: attempt to access 2-element Vector{Int64} at index [1:3]} & ``Strength 3 needs at least 3 parameters; you gave 2.'' \\
Must-run case with \code{"medium"}, which is not a listed value & \code{MethodError: Cannot `convert` an object of type Nothing to an object of type Int64} & ``must\_include case 1: \code{"medium"} is not a value of \code{level} (values: \code{"low"}, \code{"mid"}, \code{"high"}).'' \\
Rule \code{(n, x, s) -> n > 2 && x > 1.0} & \code{MethodError: no method matching isless(::Float64, ::Nothing)} & Unreachable with scoped rules (\S\ref{sec:solvers}). \\
Implicit constraint & \code{BoundsError: attempt to access 2-element Vector{Symbol} at index [0]} & ``\code{os = :windows, gpu = true} cannot appear in any valid case (rules 1 and 2 together); left out.'' \\
Same, with GND & Never returns & The same message. \\
The guide's \code{wayness} example & \code{MethodError: Cannot `convert` an object of type Int64 to an object of type UnitRange{Int64}} & Accept ranges. \\
\code{wayness} at the base strength, GND & \code{AssertionError: minimum(orders) > base_wayness} & Ignore the redundant entry, or explain it. \\
\end{xltabular}
}

% =====================================================================
\section{Communicating what it does, that it did it, and what is possible next}\label{sec:communicate}

This is the center of your question. A user passes through four stages, and
at each one they have a question the package can answer better than they
can. The design principle is simple: \textbf{answer the question at the
moment it arises, in the user's own vocabulary, and point to the next useful
step.}

\begin{table}[h]
\centering\small
\caption{The user's question at each stage, and who answers it.}\label{tab:stages}
\begin{tabularx}{\linewidth}{@{}L{0.12\linewidth} L{0.36\linewidth} Y@{}}
\toprule
Stage & The user's question & Answered today by \\
\midrule
Before & Is this the right tool? What will it cost? Which strength should I choose? & The user's own arithmetic \\
Generate & Did it keep its promise? What did it leave out, and why? Can I reproduce it? & Nothing \\
Run & Which case is this? Which cases failed, by name? & A positional \code{Vector{Any}} \\
After & What caused the failure? What remains unverified? What would it cost to go further? & Nothing \\
\bottomrule
\end{tabularx}
\end{table}

\subsection{Before: a preview that helps choose a strength}

The decision a user struggles with is the strength $t$. The package can make
it concrete, because IPOG is fast enough to simply try each option. For a CI
space of five parameters with two rules, the numbers below are real
(\emph{measured}). Only the presentation is a proposal.

\begin{term}
julia> compare_strengths(space; forbid = rules)
5 parameters (2-3 values each): 72 combinations, 48 allowed by the rules.

 strength  cases  share of allowed  pairs covered  triples covered
    1        3          6%              49%             24%
    2       10         21%             100%             69%
    3       20         42%             100%            100%
    4       39         81%             100%            100%

Most failures involve 1-2 parameters (Kuhn et al. 2004); strength 2 is
the usual start. Raise strength for risky subsets with `stronger = ...`.
\end{term}

\subsection{Generate: a result that states its guarantee}

The most valuable single change is to return an object that is still an
\code{AbstractVector}, so every existing \code{for} loop keeps working, but
that prints a summary. Here is the same CI example. Every number comes from
the real run.

\begin{term}
julia> design = all_pairs(space; forbid = [
           (os, gpu)  -> os == :mac && gpu,
           (os, arch) -> os == :windows && arch == :arm64])
TestDesign: 10 cases over 5 parameters (strength 2, IPOG, deterministic)
  Guarantee  every feasible pair of values appears in at least one case
             55 of 55 feasible pairs covered (checked)
  Space      72 combinations in full; 48 satisfy the rules
  Left out   2 pairs forbidden by rules:
               os = :mac,     gpu = true       rule 1
               os = :windows, arch = :arm64    rule 2
  Bonus      83 of 120 feasible triples also covered (69%)
  Budget     first 5 cases cover 73% of pairs; first 8 cover 95%
  Next       strength 3 needs 20 cases: extend(design; strength = 3)

   #  os        julia   threads  gpu    arch
   1  :linux    "1.10"  1        false  :x64
   2  :linux    "1.11"  4        true   :arm64
   3  :linux    "1.12"  1        true   :x64
   4  :mac      "1.10"  4        false  :arm64
   5  :mac      "1.11"  1        false  :x64
      ... 5 more
\end{term}

When the rules make a pair impossible, the ``Left out'' section says so and
says why. It distinguishes pairs a rule forbids directly from pairs that are
impossible only because rules combine. The prototype already produces this
list for the three-parameter driver example:

\begin{term}
  Left out   3 pairs cannot appear in any valid case:
               os = :windows, driver = :cuda    rule 2
               gpu = true, driver = :none       rule 1
               os = :windows, gpu = true        impossible: rules 1 and 2 together
\end{term}

This one screen answers ``did it keep its promise'', ``what did it leave
out'', ``can I stop early'', and ``what next''. The user does not have to
ask.

\subsection{Run: cases with names, outcomes that are kept}

With named parameters, each case is a \code{NamedTuple}, so test output reads
\code{(n = 1000, method = :newton, tol = 0.001, sparse = true)}. A small
helper can run the body for every case. It gives each case its own named
testset, keeps going after failures, and keeps each outcome next to the
design:

\begin{jlcode}
@testset "solver" begin
    global results = run_cases(design) do c
        @test solve(c.n, c.method, c.tol, c.sparse) > 0
    end
end
\end{jlcode}

Simulation developers do not run cases inside Test.jl, so the same idea
needs a file form. They write the design out, run the jobs, and read the
outcomes back in (\code{record_outcomes(design, "results.csv")}), so that
\code{diagnose} and \code{coverage} work on campaigns as well as unit tests.

\subsection{After: diagnosis, masking, and the next step}

This output is again computed from the real run in
Section~\ref{sec:ev-failure}:

\begin{term}
julia> diagnose(results)
2 of 10 cases failed, both with ErrorException("singular factorization").

Combinations present only in failing cases:
  method = :newton, sparse = true     2 of 2 failures, 0 of 8 passes   <- prime suspect
  n = 1000, method = :newton          1 of 2 failures
  n = 100,  sparse = true             1 of 2 failures

Not yet verified: the last two pairs appeared only in failing cases.
Two follow-up cases would settle all three:
  (n = 1000, method = :newton, tol = 0.001,  sparse = false)
  (n = 100,  method = :gmres,  tol = 1.0e-6, sparse = true)
Run them with: run_cases(followups(results)) do c ... end
\end{term}

Three more functions complete the loop after a run:

\begin{itemize}[leftmargin=1.5em]
\item \code{extend(design; strength = 3)} keeps the existing cases and adds
  only what the higher strength needs. The seed mechanism already does this
  internally.
\item \code{coverage(cases, space; strength = 2)} measures \emph{any} list of
  cases, including hand-written ones, and lists what is missing by name.
\item \code{topup(cases, space)} returns only the new cases needed to reach
  full coverage, for maintainers adopting the package gradually.
\end{itemize}

\subsection{For AI agents in particular}

Everything above helps agents, and a few further things matter especially
to them:

\begin{itemize}[leftmargin=1.5em]
\item \textbf{A machine-readable summary.} \code{summary(design)} returns a
  \code{NamedTuple} (cases, strength, engine, seed, required, covered,
  forbidden, impossible, full size) that serializes to JSON, so an agent can
  report exact figures instead of paraphrasing.
\item \textbf{Checkable claims.} \code{coverage} is independent of the
  generator, so an agent, or the human reviewing its work, can verify the
  guarantee in one line.
\item \textbf{Docstrings that begin with the situation.} For example: ``Use
  when a function or system has several discrete options and testing every
  combination is too slow.'' Agents choose tools by matching situations.
\item \textbf{Errors that name things and suggest a fix.} Agents repair code
  from messages, and a message that names \code{level} and lists its values
  gets fixed on the first try.
\item \textbf{Literal output.} \code{as_code(design)} prints the cases as a
  Julia literal, or as a \code{@testset} block, so the committed test file
  shows every case. Reviewers see exactly what runs, and the test has no
  hidden generation step.
\item \textbf{A one-page guide in the repository} (an \code{llms.txt}-style
  file) giving the decision rule and three canonical patterns: test a
  function, plan a CI matrix, audit existing tests.
\end{itemize}

\subsection{Words matter}

Some of the current names get in the way of understanding
(Table~\ref{tab:vocab}). The old names can stay as aliases.

\begin{table}[h]
\centering\small
\caption{Suggested vocabulary.}\label{tab:vocab}
\begin{tabularx}{\linewidth}{@{}L{0.25\linewidth} Y L{0.26\linewidth}@{}}
\toprule
Today & Problem & Suggested \\
\midrule
\code{seeds} & Collides with random-number seeds. The package itself uses both, and so does its test suite (\code{seed_mod}). & \code{must_include} \\
\code{n_way} & Fine for experts; the literature says ``strength''. & \code{strength} (keep \code{n_way}) \\
\code{wayness = Dict(3 => [[3,4,5,6]])} & Dictionary of lists of lists of positions. & \code{stronger = [(:a, :b, :c) => 3]} \\
\code{disallow} over all positions & Must handle \code{nothing}; the scope is invisible. & \code{forbid} rules over named arguments \\
``coverage'' & Collides with line coverage. & ``interaction coverage'' \\
\code{IPOG}, \code{GND} & Named by algorithm, not by behavior. & Keep, but document as ``deterministic'' and ``randomized, varies by seed'' \\
\bottomrule
\end{tabularx}
\end{table}

% =====================================================================
\section{A sketch of the interface}\label{sec:api}

The sketch follows four principles. It stays backward compatible: today's
positional calls keep working. It prefers names over positions. The result
\emph{is} the vector, so loops keep working. And every function returns
something that prints a useful summary.

\begin{jlcode}
using UnitTestDesign

space = TestSpace(
    os      = [:linux, :mac, :windows],
    julia   = ["1.10", "1.11", "1.12"],
    threads = [1, 4],
    gpu     = [false, true],
    arch    = [:x64, :arm64],
)

design = all_pairs(space;
    forbid       = [(os, gpu)  -> os == :mac && gpu,             # scope = argument names
                    (os, arch) -> os == :windows && arch == :arm64],
    must_include = [(os = :linux, julia = "1.12", threads = 4, gpu = true, arch = :x64)],
    stronger     = [(:julia, :threads, :gpu) => 3],
)

design[3]                         # (os = :linux, julia = "1.12", threads = 1, ...)
coverage(design)                  # checked independently of the generator
coverage(handwritten, space)      # audit an existing suite
topup(handwritten, space)         # only the cases needed to complete it
extend(design; strength = 3)      # keep these cases, add the rest

github_matrix(design)             # JSON for a workflow's `include:` list
write_design("campaign.toml", design); read_design("campaign.toml")
as_code(design)                   # literal cases for a reviewable test file
\end{jlcode}

\paragraph{Invalid values.}
A value wrapped as \code{Invalid(-1)} would be placed with only valid values
around it, and every invalid value would still be paired with every valid
value of the other parameters~\cite{pict}. The summary would report the two
guarantees separately.

\paragraph{Partitions with random representatives.}
This makes the equivalence-class idea from the paper a first-class feature.
A value can be a named partition that draws a concrete value at run time
from a recorded seed:

\begin{jlcode}
space = TestSpace(
    x   = [Partition(:negative, rng -> -1e3 * rand(rng)),
           Partition(:tiny,     rng -> 1e-12 * rand(rng)),
           0.0, Inf],
    alg = [:rk4, :midpoint],
)
\end{jlcode}

Coverage is guaranteed over partitions, while values vary from run to run.
This is the partnership with property-based testing described in
Section~\ref{sec:moment}, and a package extension could accept Supposition.jl generators
directly.

\paragraph{Budgets and rolling coverage.}
\code{all_pairs(space; max_cases = 12)} would return the best 12 cases in
coverage order and report the share covered. \code{GND(seed = run_number)} in
CI would give a different valid pairwise matrix on each run. The summary
could then report how much three-way coverage has built up across recorded
runs.

% =====================================================================
\section{Making the solvers reliable}\label{sec:solvers}

\subsection{Two ways generation fails}

Constrained generation fails in two ways, and one design fixes both.

\begin{enumerate}[leftmargin=1.5em]
\item \textbf{Implicit constraints.} A tuple passes the rule on its own, but
  no complete, allowed test contains it. The engines treat it as required.
  IPOG then crashes, and GND waits forever. This accounts for 28 of 43
  pairwise crashes and 65 of 69 three-way crashes.
\item \textbf{Greedy dead ends.} Every required tuple is feasible, but values
  are placed after checking only that the row is not \emph{yet} forbidden. A
  row can become impossible to complete. This accounts for 15 and 4
  crashes.
\end{enumerate}

\subsection{The fix: feasibility first, then keep every row completable}

\begin{proposalbox}{Four parts.}
\begin{enumerate}[leftmargin=1.3em]
\item \textbf{Feasibility pass.} A tuple is required only if some complete,
  allowed test contains it. Check each tuple once, before generation.
  Report the rest by name, and distinguish ``forbidden by rule $k$'' from
  ``impossible because rules combine''.
\item \textbf{Completability invariant.} Whenever an engine sets a value,
  require that the row can still be completed. In IPOG that means three
  sites: \code{choose_last_parameter_filter!} (line 34),
  \code{insert_tuple_into_tests_filter}, and
  \code{fill_remaining_missing_values_filter!} (line 226). In GND it means
  candidate construction.
\item \textbf{Guaranteed progress.} With parts 1 and 2, IPOG cannot reach a
  dead end. GND always has a covering candidate, because the first uncovered
  tuple is itself completable. Every remaining loop should still stop, with
  a diagnosis, if a pass makes no progress.
\item \textbf{Final validation.} Check every complete row, including
  must-run cases, and report any case that breaks a rule. Reject it unless
  the user insists.
\end{enumerate}
\end{proposalbox}

My prototype implements parts 1--3 around a simple greedy generator. It uses
a depth-first search to check completability. It returned valid, fully
covering designs for all 1,000 random problems, including every problem that
crashed IPOG, and it produced the ``impossible: rules 1 and 2 together''
report shown earlier. It is not tuned for design size. Its purpose is to
show that the completability invariant alone removes both failure classes.

\paragraph{Cost.}
The completability check is a backtracking search over unassigned
parameters. In the worst case it is exponential, but in practice it is cheap,
because rules are local. With scoped rules, only parameters connected to the
row's assigned parameters through the rule graph need to be searched.
Unconstrained parameters can be filled in freely, and results can be cached
on the constrained part of the row. This mirrors the optimizations in IPOG-C,
which reduced constraint-checking cost by one to two orders of magnitude and
ships in NIST's ACTS~\cite{yu2013}. IPOG-C is the closest published model for
this change.

\subsection{Rules the solver can understand}

Scoped rules, meaning functions whose argument names are parameter names,
fix several problems at once:

\begin{itemize}[leftmargin=1.5em]
\item The engine calls a rule only when all of its parameters are assigned,
  so \code{nothing} never appears. A real \code{nothing} value is just a
  value.
\item The scope is known, which the cheap completability check above
  depends on.
\item Messages can name the rule (``rule 2: \code{os}, \code{arch}''), and a
  label can be attached by writing a rule as a pair,
  \code{"no CUDA on Windows" => rule}.
\end{itemize}

\code{Base.method_argnames} recovers the names. I checked that it works on
Julia 1.12 for both anonymous and named functions. It is internal to Base,
however, so the package should also accept an explicit form such as
\code{(:os, :gpu) => (o, g) -> o == :mac && g}. The current whole-row
\code{disallow} can stay for compatibility. For it, the engine would check
complete rows only and search over all parameters, which is slower but
correct.

\subsection{Where Z3 fits}

The \code{z3_example.jl} draft translates one kind of expression: a parameter
compared with a literal, combined with \code{&&} and \code{||}. That is
narrower than what users can write today, which is any Julia predicate. It
also adds a large binary dependency, for domains of a handful of values where
a short search is enough. I would keep Z3 out of the core. If very large
constrained spaces appear, an optional package extension could translate
scoped rules into SMT and speed up the same completability question.

\subsection{Smaller fixes}

\begin{itemize}[leftmargin=1.5em]
\item Change \code{n_way} to \code{n_way - 1} in the GND match condition
  (\code{src/coverage_matrix.jl:193}). Give GND a fixed default seed, record
  it in the design, and print it when a test fails.
\item Copy the \code{wayness} dictionary rather than modifying it. Accept
  ranges and tuples in it. Ignore, or explain, entries at the base strength
  instead of asserting.
\item Allow single-valued parameters, and a single parameter at strength 1.
  Check that the strength does not exceed the number of parameters.
\item Validate must-run cases by name and value before generating.
\item Have \code{full_factorial} estimate its size and refuse beyond a limit
  unless asked.
\item Say when an excursion's base case is forbidden.
\item Replace column slices in hot loops with views or explicit loops, which
  would remove most of the multi-gigabyte allocation.
\item Add the random constrained problems and the independent checker to the
  test suite. They would have caught both failure classes and the GND bug.
\end{itemize}

% =====================================================================
\section{Documentation and first contact}

\begin{itemize}[leftmargin=1.5em]
\item \textbf{The README's first screen} should state the situation (``many
  options; too many combinations to test all''), the guarantee in one
  sentence, one example with named parameters and its printed summary, and a
  line on when to use something else.
\item \textbf{Organize the manual by job, not by algorithm.} Useful recipes:
  test a function with many options; test generic code across types; plan a
  CI matrix; run a simulation campaign from files; audit an existing suite;
  diagnose a failure; test invalid inputs; combine with property-based
  testing. Algorithm pages can stay, further down.
\item \textbf{Move the advice on choosing values into the manual.}
  Equivalence classes, boundary values, and how many values to list are the
  most important human decisions in combinatorial testing, and today they
  are explained well only in the paper.
\item \textbf{Correct the engine page.} Make every example a doctest; that
  alone would have caught the broken \code{wayness} example.
\item \textbf{Add the one-page agent guide} described in
  Section~\ref{sec:communicate}.
\end{itemize}

% =====================================================================
\section{Roadmap}

\begin{table}[h]
\centering\small
\caption{A phased plan. Each phase is useful by itself.}\label{tab:roadmap}
\begin{tabularx}{\linewidth}{@{}L{0.17\linewidth} Y L{0.12\linewidth}@{}}
\toprule
Phase & Work & Scale \\
\midrule
0. Correct & GND one-token fix and default seed; feasibility pass and completability checks; no loops without progress; must-run case validation; \code{wayness} fixes; user-vocabulary errors; engine documentation; random-rule regression tests. & 1--2 weeks \\
1. Speak up & Design object with a printed summary; named parameters and \code{NamedTuple} cases; scoped rules; \code{coverage}; \code{must_include} and \code{strength} names, with aliases. & 2--4 weeks \\
2. Close the loop & \code{run_cases} and file-based outcomes; \code{diagnose} and \code{followups}; \code{extend} and \code{topup}; \code{max_cases} and coverage ordering. & 3--4 weeks \\
3. Reach & Tables.jl extension; CSV/TOML; GitHub matrix; \code{as_code}; \code{Invalid}; \code{Partition}; rolling coverage; agent guide; job-oriented recipes. & Ongoing \\
\bottomrule
\end{tabularx}
\end{table}

Phase 0 is the one I would not postpone. A user with rules is exactly the
user who needs this package most, and in my random constrained problems about
one in nine ended in a confusing crash; GND users get a process that never ends.

% =====================================================================
\section{Closing}

The hard part of this package, the covering algorithms, is already done and
largely correct. It is fast, it has few dependencies, and its documentation
and paper explain the ideas generously. What it lacks is a voice while it
runs. It never tells users what it promised, what it left out, what
happened, or what they could do next.

The users this package could serve best are people who already know that
some combinations matter and cannot afford all of them: scientists with
expensive runs, maintainers of generic code, engineers paying for CI, and
agents writing tests for them. Each of them would trust the tool more if it
showed its work, and would get more from it if it helped with the step after
generation. Named parameters, a summary that states the guarantee, reliable
rules, and a diagnosis after failure would make \utd the tool people reach
for at the moment the loop explodes.

% =====================================================================
\appendix
\section{Reproducing the experiments}\label{app:repro}

All scripts ran against the working tree (v0.4.0, commit \code{e0504d6}),
with the package added by \code{Pkg.develop} to a scratch environment. The
experiments changed no repository files, and the scripts live outside the
repository; the constraint and GND scripts are the ones worth turning into
regression tests. The checker in \code{verify.jl} enumerates the
full factorial to find feasible tuples, so it is limited to small problems
by design.

\begin{center}\small
\begin{tabularx}{\linewidth}{@{}L{0.14\linewidth} Y@{}}
\toprule
Script & What it measures \\
\midrule
\code{e1.jl} & First contact: REPL output, exported names, edge-case errors \\
\code{e2.jl}, \code{e3.jl} & 1,000 random constrained problems; classification of crashes \\
\code{e4.jl} & GND on the implicit-constraint example under a 45-second watchdog \\
\code{e5.jl} & \code{wayness} example, dictionary mutation, GND reproducibility, IPOG stability \\
\code{e6.jl} & Mixed-strength verification; triples covered by pairwise designs; coverage by prefix \\
\code{e7.jl}, \code{e8.jl} & Time, allocation, and size for IPOG and GND; GND sensitivity to \code{M} \\
\code{e9.jl} & GND sizes before and after the one-token patch (patched in-session only) \\
\code{e10.jl}, \code{e16.jl} & Test.jl output for a failing design; localization and masked pairs \\
\code{e11.jl} & Excursion base case, rule-violating must-run cases, value types \\
\code{e12.jl}, \code{proto.jl} & Constraint-safe prototype on the same 1,000 problems \\
\code{e14.jl}, \code{e17.jl} & Figures for the CI example and the strength comparison \\
\code{e15.jl} & Silent coverage loss when \code{nothing} is a real value \\
\bottomrule
\end{tabularx}
\end{center}

\section{Findings by location}

\begin{center}\small
\begin{tabularx}{\linewidth}{@{}L{0.36\linewidth} Y@{}}
\toprule
Location & Finding \\
\midrule
\code{src/coverage_matrix.jl:193} & GND scoring uses \code{n_way}; needs \code{n_way - 1}. \\
\code{src/greedy_tuples.jl:195, 201} & Loops with no progress check; hang on uncoverable tuples or all-forbidden candidates. \\
\code{src/greedy_tuples.jl:250, 256} & The same loops in the mixed-strength variant. \\
\code{src/parameter_order.jl:128, 170} & Unplaceable tuples appended as rows without a completability check. \\
\code{src/parameter_order.jl:224--226} & Final fill leaves zeros when no value is allowed; leads to \code{p[0]}. \\
\code{src/parameter_order.jl:436}, \code{src/excursions.jl:32} & The user's \code{wayness} dictionary is modified. \\
\code{src/coverage_matrix.jl:426} & Assertion instead of a message when an extra strength equals the base. \\
\code{src/factorial_interface.jl:61} & GND's default RNG is unseeded and unrecorded. \\
\code{src/factorial_interface.jl:85} & \code{nothing} marks unassigned parameters; ambiguous with real \code{nothing}. \\
\code{src/factorial_interface.jl:100} & Must-run values looked up with \code{indexin}; a miss surfaces as a \code{convert} error. \\
\code{src/factorial_interface.jl:219} & Single-valued parameters rejected. \\
\code{src/full_factorial.jl:5} & Allocates the full result before any size check. \\
\code{docs/src/man/guide.md} & \code{wayness} example throws. \\
\code{docs/src/man/engines.md} & Describes GND as shorter; measured longer. \\
\bottomrule
\end{tabularx}
\end{center}

% =====================================================================
\begin{thebibliography}{99}\small

\bibitem{kuhn2004}
D.~R. Kuhn, D.~R. Wallace, and A.~M. Gallo Jr.
\newblock Software fault interactions and implications for software testing.
\newblock \emph{IEEE Transactions on Software Engineering}, 30(6):418--421, 2004.

\bibitem{lei2008}
Y.~Lei, R.~Kacker, D.~R. Kuhn, V.~Okun, and J.~Lawrence.
\newblock IPOG/IPOG-D: Efficient test generation for multi-way combinatorial testing.
\newblock \emph{Software Testing, Verification and Reliability}, 18(3):125--148, 2008.

\bibitem{yu2013}
L.~Yu, Y.~Lei, M.~Nourozborazjany, R.~N. Kacker, and D.~R. Kuhn.
\newblock An efficient algorithm for constraint handling in combinatorial test generation.
\newblock In \emph{IEEE Sixth International Conference on Software Testing, Verification and Validation (ICST)}, 2013.

\bibitem{ccm}
D.~R. Kuhn, I.~Dominguez Mendoza, R.~N. Kacker, and Y.~Lei.
\newblock Combinatorial coverage measurement concepts and applications.
\newblock In \emph{Second International Workshop on Combinatorial Testing (IWCT)}, 2013.
\newblock Tool: \url{https://github.com/usnistgov/combinatorial-testing-tools}.

\bibitem{ben}
L.~S. Ghandehari et al.
\newblock A combinatorial testing-based approach to fault localization.
\newblock \emph{IEEE Transactions on Software Engineering}, 46(6):616--645, 2020.

\bibitem{yilmaz2014}
C.~Yilmaz, E.~Dumlu, M.~B. Cohen, and A.~Porter.
\newblock Reducing masking effects in combinatorial interaction testing: A feedback driven adaptive approach.
\newblock \emph{IEEE Transactions on Software Engineering}, 40(1):43--66, 2014.

\bibitem{pict}
Microsoft.
\newblock PICT: Pairwise Independent Combinatorial Testing, user guide.
\newblock \url{https://github.com/microsoft/pict/blob/main/doc/pict.md}.

\bibitem{supposition}
Seelengrab (GitHub author).
\newblock Supposition.jl: choice-sequence based property-based testing for Julia.
\newblock \url{https://github.com/Seelengrab/Supposition.jl}.

\bibitem{ghmatrix}
GitHub Docs.
\newblock Running variations of jobs in a workflow (a matrix generates at most 256 jobs per workflow run).
\newblock \url{https://docs.github.com/en/actions/using-jobs/using-a-matrix-for-your-jobs}.

\end{thebibliography}

\end{document}
